\documentclass[aspectratio=169]{beamer}
\usetheme{Madrid}
\usecolortheme{seahorse}
\usepackage{graphicx}
\usepackage{booktabs}
\usepackage{xcolor}
\graphicspath{{figs/}}
\setbeamertemplate{navigation symbols}{}
\setbeamertemplate{caption}{\raggedright\insertcaption\par}
\newcommand{\fig}[2]{\IfFileExists{figs/#1}{\includegraphics[width=#2]{#1}}{\fbox{\textit{[figure: #1]}}}}
\newcommand{\good}[1]{\textcolor{green!55!black}{\textbf{#1}}}
\newcommand{\bad}[1]{\textcolor{red!70!black}{\textbf{#1}}}

\title[Fuzzy Soil Index — Replication]{Does a ``Fuzzy Soil Index'' Help a Robot Walk on Soft Ground?}
\subtitle{A replication study of the \textit{wcci2026\_hunter} fuzzy soil model}
\author{Hunter humanoid \,/\, IsaacSim \,/\, HumanoidVerse}
\date{\today}

\begin{document}

\frame{\titlepage}

\begin{frame}{The problem (from the paper)}
  \textbf{Goal:} make a humanoid robot walk on \textbf{agricultural soil} — furrowed,
  deformable ground where the feet sink and slip.
  \begin{itemize}
    \item Real soil is \textbf{not uniform}: it changes from firm \& grippy to soft \& slippery.
    \item Two things vary independently:
      \begin{itemize}
        \item \textbf{Traction} — how grippy the surface is (friction $\mu$).
        \item \textbf{Support} — how firm it is / how much the foot sinks (stiffness).
      \end{itemize}
    \item The robot must \textbf{judge how hard the current soil is} to adapt and stay upright.
  \end{itemize}
  \vfill
  \begin{block}{The paper's pitch}
    Combine traction + support into one human-readable ``\textbf{soil difficulty score}''
    using \textbf{fuzzy logic}, and use it to drive training / reduce foot slip ($>$95\% claimed).
  \end{block}
\end{frame}

\begin{frame}{What the ``fuzzy soil index'' actually is}
  A \textbf{fuzzy (Mamdani) controller}: two inputs $\to$ one difficulty score.
  \begin{columns}
    \begin{column}{0.55\textwidth}
      \begin{itemize}
        \item Inputs: \textbf{traction} $\mu$ and \textbf{support} (stiffness).
        \item Each is graded \textit{Low / Med / High} with \textbf{smooth} (fuzzy) boundaries.
        \item Rules $\to$ soil is \textit{Easy / Moderate / Hard}.
        \item Output: one number $d\in[0.2,0.8]$ (higher = harder).
      \end{itemize}
      \textbf{Selling points:} interpretable bins, smooth interpolation, blends two factors.
    \end{column}
    \begin{column}{0.42\textwidth}
      \centering
      \fbox{\parbox{0.95\linewidth}{\centering
        \textbf{Two soil ``dials''}\\[4pt]
        traction $\mu$ \;\&\; support (firmness)\\[6pt]
        $\downarrow$ fuzzy rules $\downarrow$\\[6pt]
        \textbf{difficulty} = Easy / Mod / Hard}}
    \end{column}
  \end{columns}
\end{frame}

\begin{frame}{Our question}
  \begin{block}{}
    \centering \Large Does the fuzzy index \textbf{actually add value} over just using the
    \textbf{raw soil numbers}?
  \end{block}
  \vfill
  \begin{itemize}
    \item A fuzzy score is only useful if it \textbf{beats the raw measurements} it is built from.
    \item We replicate the recipe on the \textbf{Hunter} humanoid in \textbf{IsaacSim} (HumanoidVerse)
          and test the index four ways.
    \item We judge it as a \textbf{difficulty descriptor} and as a \textbf{control signal}.
  \end{itemize}
  \vfill
  \footnotesize Honest-negative friendly: we want the right answer, not a flattering one.
\end{frame}

\begin{frame}{How we tested it — four settings}
  \centering
  \renewcommand{\arraystretch}{1.25}
  \begin{tabular}{@{}llll@{}}
    \toprule
    \textbf{Study} & \textbf{Simulator} & \textbf{Fuzzy used as} & \textbf{Support axis} \\
    \midrule
    A   & rigid PhysX        & passive \textbf{label}       & physically \bad{inert} \\
    B   & rigid PhysX        & \textbf{curriculum} controller & inert \\
    C   & \textbf{DFH deformable} & \textbf{descriptor}      & physically \good{active} \\
    \,\textcircled{1}\,\textcircled{2} & DFH (+ noisy sensing) & descriptor / decisions & active \\
    \bottomrule
  \end{tabular}
  \vfill
  \footnotesize Each study closes the previous one's loophole. ``Support axis'' = whether soil
  \textit{firmness} actually does anything in that simulator.
\end{frame}

\begin{frame}{Sets A \& B (rigid simulator)}
  \textbf{Set A — fuzzy as a label:} does $d$ track difficulty / beat raw $\mu$?
  \begin{itemize}
    \item $d$ orders soil correctly (Spearman $\rho\approx0.97$) \dots
    \item \dots but adds \bad{nothing} over raw friction $\mu$.
    \item \textbf{Why:} rigid PhysX has \textbf{no real ``firmness'' axis} — the foot can't sink,
          so the 2-dial model collapses to 1 dial (friction).
  \end{itemize}
  \vspace{4pt}
  \textbf{Set B — fuzzy drives the training curriculum:}
  \begin{itemize}
    \item fuzzy pacing $\approx$ crisp threshold $\approx$ fixed schedule (no significant difference).
  \end{itemize}
  \vfill
  \begin{block}{Catch}
    A 2-input fuzzy model can't shine when \textbf{one input does nothing}.
    \textbf{We need a simulator where soil firmness is physically real.}
  \end{block}
\end{frame}

\begin{frame}{The fix: DFH deformable soil}
  \textbf{DFH} (Deformable Furrowed Heightfield) adds the missing physics:
  feet \textbf{sink} (Bekker), generate \textbf{sideways drag}, and the soil \textbf{deforms}.
  \begin{columns}
    \begin{column}{0.5\textwidth}
      \centering \fig{dfh_support_axis.png}{\linewidth}
    \end{column}
    \begin{column}{0.48\textwidth}
      \good{The support axis is now real:}
      \begin{itemize}
        \item Soft soil ($K{=}1$): feet sink $\sim$\textbf{49\,mm}.
        \item Firm soil ($K{=}8$): sink $\sim$\textbf{16\,mm}; drag drops $33\to9$\,N.
        \item Confirmed by a force-off control (no forces $\Rightarrow$ no effect).
      \end{itemize}
      \textit{Same robot, same gait — only the soil firmness changes. You can see it.}
    \end{column}
  \end{columns}
\end{frame}

\begin{frame}{Set C: the key result}
  Now \textbf{both} soil dials are physical. Does the fuzzy index finally add value?
  \begin{columns}
    \begin{column}{0.52\textwidth}
      \centering \fig{dfh_fuzzy_vs_raw.png}{\linewidth}
    \end{column}
    \begin{column}{0.46\textwidth}
      \begin{itemize}
        \item Raw soil numbers \good{strongly predict} difficulty (+0.33 to +0.42).
        \item The paper's \textbf{fuzzy index} adds \bad{+0.025} — essentially nothing
              (permutation $p{=}0.001$).
      \end{itemize}
      \begin{block}{Sharp finding}
        Even where firmness \textbf{is} physically active, the fuzzy mapping \textbf{fails to use it}.
        The problem is the \textbf{fuzzy model's structure}, \textbf{not} the simulator.
      \end{block}
    \end{column}
  \end{columns}
\end{frame}

\begin{frame}{Giving fuzzy its best shot: noisy soil sensing}
  Fuzzy logic's real strength is \textbf{tolerating imprecise inputs}. A real robot only
  \textit{estimates} the soil. So we added sensing noise.
  \begin{columns}
    \begin{column}{0.55\textwidth}
      \centering \fig{fuzzy_noise_analysis.png}{\linewidth}
    \end{column}
    \begin{column}{0.43\textwidth}
      \begin{itemize}
        \item Fuzzy \textbf{is} noise-stable (flat line) and overtakes raw \textit{only} at
              extreme noise.
        \item But it wins by being \bad{stable, not accurate} — raw still \textbf{orders soils better}.
        \item \textbf{Decisions} (\textcircled{2}): a fuzzy-driven curriculum \bad{camps on easy soil}
              and never advances.
      \end{itemize}
    \end{column}
  \end{columns}
  \vfill
  \footnotesize Stability of an \textit{uninformative} signal — it doesn't swing because it barely says anything.
\end{frame}

\begin{frame}{Verdict}
  \begin{block}{Bottom line}
    Across rigid \textbf{and} deformable soil, the paper's fuzzy soil index shows
    \textbf{no demonstrable value} over the raw soil numbers — as a label, a curriculum,
    a descriptor, or under noisy sensing.
  \end{block}
  \begin{itemize}
    \item \good{Replicates:} the friction-curriculum \textbf{survival} benefit; DFH genuinely
          restores the firmness axis.
    \item \bad{Does not add value:} the fuzzy mapping collapses two soil factors into a coarse
          score and \textbf{discards the useful information}.
    \item The failure is \textbf{structural}, not a simulator artifact (Set C proves it).
  \end{itemize}
\end{frame}

\begin{frame}{Scope \& what would change the answer}
  \textbf{Honest scope:}
  \begin{itemize}
    \item This is about \textbf{this paper's specific fuzzy mapping} — \textit{not} all 2-D or
          learned terrain descriptors.
    \item DFH is an \textbf{uncalibrated} deformable-soil surrogate; one frozen walker.
    \item Not a disproof of the paper on \textbf{real} soil — a non-replication \textit{in simulation}.
  \end{itemize}
  \textbf{What could make ``fuzzy matter'':}
  \begin{itemize}
    \item Redesign the fuzzy rules so they \textbf{actually use the firmness axis} (they ignore it
          at mid-traction), or use a \textbf{learned} 2-D descriptor.
    \item Train a soil-robust walker (removes the single-policy caveat).
  \end{itemize}
\end{frame}

\begin{frame}{Summary}
  \centering
  \Large \textbf{A fuzzy score is only useful if it beats the numbers it's made from.}\\[10pt]
  \normalsize Here, it doesn't.\\[16pt]
  \footnotesize
  Code \& full write-ups: \texttt{docs/experiments/fuzzy\_soil\_UNIFIED\_conclusion.md};
  models on Hugging Face \texttt{anhrisn/hunter-dfh-locomotion}.
\end{frame}

\end{document}
